\section*{第 \ref{ch:b1:derivative} 章——求导}

\begin{solution}{exo:b1:derivative:1}
$x^x = \eu^{x\ln x}$ 在 $\intoo{0}{+\infty}$ 上，导数为 $(\ln x +
1)\,x^x$。

$\ln(x + \sqrt{x^2+1})$ 在 $\R$ 上（真数恒 $> 0$）：
导数为 $\frac{1}{\sqrt{x^2+1}}$（在
\cref{prop:b1:functions:invhyp}中已算过，此即
$\operatorname{arsinh}$）。

$\arctan\frac1x$ 在 $\R^*$ 上，导数为 $\frac{-1/x^2}{1 + 1/x^2} =
\frac{-1}{1 + x^2}$（与
\cref{prop:b1:functions:arcidentities}~(2)一致：函数在
每条半直线上为 $\pm\frac\pi2 - \arctan x$）。

$\sqrt{1 + \eu^{2x}}$ 在 $\R$ 上，导数为
$\frac{\eu^{2x}}{\sqrt{1 + \eu^{2x}}}$。
\end{solution}

\begin{solution}{exo:b1:derivative:2}
在 $0$ 处：$\bigl|\frac{f(h) - 0}{h}\bigr| = \abs{h \sin\frac1h} \leq
\abs h \to 0$，故 $f'(0) = 0$。当 $x \neq 0$ 时，由通常的求导法则可得
$f'(x) = 2x\sin\frac1x - \cos\frac1x$。取 $x_n = \frac{1}{2\pi
n}$：$f'(x_n) = 0 - 1 \to -1$；取 $y_n = \frac{1}{(2n+1)\pi}$：
$f'(y_n) = 0 + 1 \to 1$。两个趋于 $0$ 的数列，
其对应的 $f'$ 极限不同，故无极限（见\cref{thm:b1:continuity:seqchar}），因此 $f'$
在 $0$ 处不\omterm{def:b1:continuity:continuous}{连续}，从而 $f$ \omterm{def:b1:derivative:def}{可导}但不是
$C^1$。
\end{solution}

\begin{solution}{exo:b1:derivative:3}
$\ln(1+x) < x$（$x > 0$）：由 $0$ 处的切线不等式（凹性）
（除切点外严格成立，因 $\ln$ 严格凹；
或用中值定理：$\ln(1+x) = \frac{x}{1+c}$，对某个
$c \in \intoo{0}{x}$，且 $\frac{x}{1+c} < x$）。同一
中值恒等式给出下界：$\frac{x}{1+c} > \frac{x}{1+x}$。

由此，以 $x/n$ 代替 $x$，
\[
\frac{x/n}{1 + x/n} < \ln\Bigl(1 + \frac xn\Bigr) < \frac xn
\quad\implies\quad
\frac{x}{1 + x/n} < n \ln\Bigl(1 + \frac xn\Bigr) < x .
\]
左端趋于 $x$，由夹逼原理得 $n\ln(1 + \frac xn)
\to x$，又由 $\exp$ 的\omterm{def:b1:continuity:continuous}{连续}性，$\bigl(1 + \frac xn\bigr)^n =
\eu^{n\ln(1 + x/n)} \to \eu^x$。
\end{solution}

\begin{solution}{exo:b1:derivative:4}
设 $x_1 < x_2 < \dots < x_k$ 为 $P$ 的互异根。在每个
$\intcc{x_i}{x_{i+1}}$ 上，由罗尔定理（\cref{thm:b1:derivative:rolle}）
得 $c_i \in \intoo{x_i}{x_{i+1}}$ 使 $P'(c_i) = 0$，即
$k - 1$ 个 $P'$ 的根，它们互异，原因是这些开\omterm{prop:b1:reals:intervals}{区间}
两两不交——且由构造它们彼此交错。

若 $P$（次数为 $n$）的根全为实数，记其\omterm{def:b1:poly:derivative}{重数}分别为
$m_1 + \dots + m_k = n$。\omterm{def:b1:poly:derivative}{重数} $m_i \geq 2$ 的每个根
是 $P'$ 的根，其\omterm{def:b1:poly:derivative}{重数}为 $m_i - 1$
（\cref{prop:b1:poly:multiplicity}），贡献 $\sum (m_i - 1) = n
- k$；罗尔定理再给出 $k - 1$ 个根，且都与上述根不同。合计
$\geq n - 1 = \deg P'$，故 $P'$ 的根全为实数。
\end{solution}

\begin{solution}{exo:b1:derivative:5}
取 $\varepsilon > 0$ 及 $A$，使得 $\abs{f'(t) - \ell} \leq
\varepsilon$（当 $t \geq A$ 时）。对 $x > A$，
在 $\intcc{A}{x}$ 上由中值定理得 $c \in \intoo{A}{x}$，使
\[
f(x) = f(A) + f'(c)(x - A),
\qquad\text{于是}\qquad
\Bigl|\frac{f(x)}{x} - \ell\Bigr|
\leq \frac{\abs{f(A)} + \abs\ell A}{x} + \abs{f'(c) - \ell}
\cdot\frac{x - A}{x} \leq \frac{C_A}{x} + \varepsilon .
\]
当 $x$ 充分大时，$\frac{C_A}{x} \leq \varepsilon$，因此
$\frac{f(x)}{x} \to \ell$。

是的：$f(x+1) - f(x) = f'(c_x)$，其中 $c_x \in \intoo{x}{x+1}$（
对 $\intcc{x}{x+1}$ 用中值定理），且 $c_x \to +\infty$，故
$f(x+1) - f(x) \to \ell$。
\end{solution}

\begin{solution}{exo:b1:derivative:6}
对 $n$ 用归纳法。当 $n = 1$ 时，由罗尔定理即得。若结论对 $n - 1$ 成立：
若 $f$ 在 $n+1$ 个点上为零，则在 $n$ 个间隔上用罗尔定理，
$f'$ 在 $n$ 个互异点处为零；对
$f'$ 应用归纳假设（$n-1$ 次\omterm{def:b1:derivative:def}{可导}，$n$ 个零点）后，使
$(f')^{(n-1)} = f^{(n)}$ 在某处为零。

应用：若 $P$ 为次数 $\leq n$ 的\omterm{def:b1:poly:def}{多项式}，且在 $n+1$ 个点处为零，则
$P^{(n)}$ 是一个常数，它等于 $n!$ 乘以首项系数，
因而为零，于是首项系数为 $0$，再
用向下归纳法（或直接看出所有系数均为零）得出结论。
\end{solution}

\begin{solution}{exo:b1:derivative:7}
分别在 $\intcc{a}{x_0}$ 与
$\intcc{x_0}{b}$ 上应用中值定理：
\[
f'(c_1) = \frac{f(x_0) - f(a)}{x_0 - a} = \frac{f(x_0)}{x_0 - a} > 0,
\qquad
f'(c_2) = \frac{f(b) - f(x_0)}{b - x_0} = \frac{-f(x_0)}{b - x_0} < 0,
\]
其中 $c_1 < x_0 < c_2$。再对 $f'$
在 $\intcc{c_1}{c_2}$ 上应用中值定理，得 $c$ 使得
\[
f''(c) = \frac{f'(c_2) - f'(c_1)}{c_2 - c_1} < 0 . \qedhere
\]
\end{solution}

\begin{solution}{exo:b1:derivative:8}
$f'(x) = \frac{1 - \ln x}{x^2}$：$f$ 在 $\intoc{0}{\eu}$ 上递增，
在 $\intco{\eu}{+\infty}$ 上递减，最大值 $f(\eu) =
\frac1\eu$；极限分别为 $-\infty$（于 $0^+$ 处）与 $0$（于 $+\infty$ 处）（\omterm{prop:b1:functions:powerrules}{增长
比较}）。

设 $\eu \leq a < b$，$f$ 在此\omterm{prop:b1:reals:intervals}{区间}上严格递减，于是
$\frac{\ln a}{a} > \frac{\ln b}{b}$，即 $b \ln a > a \ln b$，
亦即 $a^b > b^a$。

取 $a = \eu < b = \pi$，得 $\eu^\pi > \pi^\eu$。

较小的数对：$(2, 3)$：$2^3 = 8 < 9 = 3^2$——不等号方向相反！原因在于：
由于 $2 < \eu$，而在 $\intoo{0}{\eu}$ 上函数 $f$ 为
\emph{递增}，故两数都位于
$\eu$ 之下时比较结果翻转，而跨越 $\eu$ 时则无法一概而论（$f(2) = f(4)$
可解释 $2^4 = 4^2 = 16$ 的相等）。
\end{solution}

\begin{solution}{exo:b1:derivative:9}
\emph{对 $\ln$ 的琴生不等式，用对 $n$ 的归纳法。}断言：对正数
$x_i$ 与权重 $\lambda_i > 0$，满足 $\sum \lambda_i = 1$，则
$\ln\bigl(\sum \lambda_i x_i\bigr) \geq \sum \lambda_i \ln x_i$。当
$n = 2$ 时即为凹性。归纳步：记 $\Lambda = \lambda_1 + \dots +
\lambda_{n-1} = 1 - \lambda_n$，并令 $y = \sum_{i<n}
\frac{\lambda_i}{\Lambda} x_i$，
\[
\ln\Bigl(\sum_{i \leq n} \lambda_i x_i\Bigr)
= \ln\bigl(\Lambda y + \lambda_n x_n\bigr)
\geq \Lambda \ln y + \lambda_n \ln x_n
\geq \Lambda \sum_{i<n} \frac{\lambda_i}{\Lambda}\ln x_i
+ \lambda_n \ln x_n,
\]
先用凹性（$n = 2$），再用归纳假设。

取 $\lambda_i = \frac 1n$，$x_i = a_i$，则 $\ln\frac{\sum a_i}{n}
\geq \frac 1n \sum \ln a_i = \ln\sqrt[n]{a_1\cdots a_n}$；
两边取指数即得。等号条件：$\ln$ 是\emph{严格}凹函数（$\ln'' < 0$），
故每一步取等都迫使平均点重合——
即所有 $a_i$ 相等；反之若全都相等，则等号显然成立。
\end{solution}

\begin{solution}{exo:b1:derivative:10}
令 $g(x) = f(x) - vx$，它\omterm{def:b1:derivative:def}{可导}，且 $g'(a) = f'(a) - v < 0$，
$g'(b) = f'(b) - v > 0$。由最值定理，$g$
在 $\intcc{a}{b}$ 上于某点 $c$ 取得最小值。此点不可能是 $a$：
因 $g'(a) < 0$，$a$ 右侧附近的点满足 $g < g(a)$。此点也不可能是
$b$：因 $g'(b) > 0$，$b$ 左侧附近的点满足 $g <
g(b)$。故 $c$ 为内点，由\cref{prop:b1:derivative:fermat}
可得 $g'(c) = 0$，即 $f'(c) = v$。
\end{solution}

\begin{solution}{exo:b1:derivative:11}
\emph{唯一性：}若有两个不动点 $\ell \neq \ell'$，则
$\abs{\ell - \ell'} = \abs{f(\ell) - f(\ell')} \leq k\abs{\ell -
\ell'} < \abs{\ell - \ell'}$，矛盾。

\emph{存在性：}记 $g(x) = f(x) - x$，由中值不等式
得 $f(x) \leq f(0) + k\abs x$；于是当 $x \geq
\frac{\abs{f(0)}}{1 - k}$ 时，$g(x) \leq f(0) + kx - x \leq 0$，而
对称地，$g(-x) \geq 0$ 对充分大的 $x$ 成立。由
介值定理得 $\ell$ 为 $g$ 的零点，即为不动点。

\emph{收敛性：}再次运用中值不等式：
\[
\abs{u_{n+1} - \ell} = \abs{f(u_n) - f(\ell)} \leq k\abs{u_n - \ell},
\]
由归纳法，$\abs{u_n - \ell} \leq k^n \abs{u_0 - \ell} \to 0$。
\end{solution}

\begin{solution}{exo:b1:derivative:12}
\begin{enumerate}
  \item 若 $g(b) = g(a)$，则罗尔定理给出 $g'$ 的\omterm{def:b1:topology:closure}{内部}零点，
        此情形被排除。取 $\lambda = \frac{f(b) - f(a)}{g(b) -
        g(a)}$，并令 $h = f - \lambda g$，则 $h$ 在
        $\intcc{a}{b}$ 上\omterm{def:b1:continuity:continuous}{连续}、\omterm{def:b1:topology:closure}{内部}\omterm{def:b1:derivative:def}{可导}，且 $h(b) - h(a)
        = f(b) - f(a) - \lambda(g(b) - g(a)) = 0$。由罗尔定理
        给出 $c$，使 $h'(c) = 0$，即 $f'(c) =
        \lambda\,g'(c)$；两边除以 $g'(c) \neq 0$。
  \item 当 $x > a$ 且接近 $a$ 时，在 $\intcc{a}{x}$ 上应用 (1)
        （其中 $g' \neq 0$），得 $g(x) \neq 0$，且存在 $c_x \in
        \intoo{a}{x}$，使
        \[
        \frac{f(x)}{g(x)} = \frac{f(x) - f(a)}{g(x) - g(a)}
        = \frac{f'(c_x)}{g'(c_x)} .
        \]
        当 $x \to a^+$ 时，由夹逼知 $c_x \to a^+$，故右端
        趋于 $\ell$，即 $\frac{f}{g} \to \ell$。
  \item $\frac{f(x)}{g(x)} = x\sin\frac1x \to 0$，而
        $\frac{f'(x)}{g'(x)} = 2x\sin\frac1x - \cos\frac1x$ 在
        $0$ 处无极限（见\cref{exo:b1:derivative:2}）：
        洛必达法则只能由
        $\frac{f'}{g'}$ 的极限推出 $\frac fg$ 的极限，反之不成立。
\end{enumerate}
\end{solution}

\begin{solution}{pb:b1:derivative:1}
\textbf{1.} $\bigl|\frac ab - \frac pq\bigr| = \frac{\abs{aq -
bp}}{bq}$，且两分数不等时，$aq - bp$ 为非零整数，
其距离 $\geq \frac{1}{bq}$。因此，没有有理数
除 $x$ 本身外，进入半径为
$\frac{1}{bq}$、以 $x = \frac ab$ 为中心的去心\omterm{prop:b1:reals:intervals}{区间}。

\textbf{2.} 若 $\bigl|\sqrt2 - \frac pq\bigr| \geq 1 \geq
\frac{1}{4q^2}$，则结论成立。否则 $\frac pq \in \intoo{\sqrt2 -
1}{\sqrt2 + 1}$，故 $0 < \sqrt2 + \frac pq < 2\sqrt2 + 1 < 4$。
由 $\sqrt 2 \notin \Q$ 知 $p^2 - 2q^2$ 为非零整数，且
\[
\Bigl|\sqrt2 - \frac pq\Bigr|
= \frac{\abs{2q^2 - p^2}}{q^2\,\bigl(\sqrt2 + \frac pq\bigr)}
\geq \frac{1}{4q^2} .
\]

\textbf{3.} $(p + 2q)^2 - 2(p + q)^2 = -(p^2 - 2q^2)$：其值
$\pm1$ 沿递推传递。从 $(1,1)$ 出发：
\[
(3,2),\ (7,5),\ (17,12),\ (41,29),\ (99,70),
\]
其中 $p^2 - 2q^2$ 的值依次交替为 $-1, +1, \dots$。
对这些值，$\frac pq \geq 1$，故 $\sqrt2 + \frac pq > 2$，且
\[
\Bigl|\sqrt2 - \frac pq\Bigr| =
\frac{1}{q^2(\sqrt2 + p/q)} < \frac{1}{2q^2} ,
\]
令 $q \to \infty$，即得无穷多个 $2$ 阶逼近。
结合第 2 题，指数 $2$ 对 $\sqrt 2$ 是精确的。

\textbf{4.} 这 $N + 1$ 个数 $kx - \lfloor kx\rfloor$（$0
\leq k \leq N$）落在 $N$ 个\omterm{prop:b1:reals:intervals}{区间} $\intco{\frac
jN}{\frac{j+1}{N}}$ 中，必有两数同属一个\omterm{prop:b1:reals:intervals}{区间}
（\cref{cor:b1:counting:pigeonhole}），不妨设 $i < j$。取 $q =
j - i \leq N$，$p = \lfloor jx\rfloor - \lfloor ix\rfloor$：
$\abs{qx - p} < \frac1N$，故 $\bigl|x - \frac pq\bigr| <
\frac{1}{Nq} \leq \frac{1}{q^2}$。令 $N \to \infty$：由于
$x$ 为无理数，每个固定分数与 $x$ 的距离均为正，
而 $\frac{1}{Nq} \leq \frac 1N \to 0$，迫使
不断出现新的分数：存在无穷多个互不相同的 $\frac pq$，满足
$\bigl|x - \frac pq\bigr| < \frac{1}{q^2}$。

\textbf{5.} 任取非零整系数\omterm{def:b1:poly:def}{多项式} $P_0$，满足 $P_0(x) =
0$。若 $P_0$ 有有理根 $\frac ab$，则由因式定理
（\cref{thm:b1:poly:factor}）可写 $P_0 = \bigl(X - \frac
ab\bigr)Q$，其中 $Q \in \Q[X]$；
由于 $x \neq \frac ab$（$x$ 为无理数），$Q(x) = 0$，且
去分母后得到一个非零的\emph{整系数}\omterm{def:b1:poly:def}{多项式}，
其次数更低且以 $x$ 为根。每一步次数都会下降，
故该过程必然终止：得到 $P \in \Z[X]$，$P(x) = 0$，
且无有理根，次数为某个 $d$。若 $d \leq 1$，则 $P =
uX + v$ 将使 $x = -\frac vu$ 为有理数，故 $d \geq 2$。

\textbf{6.} $q^d\,P\bigl(\frac pq\bigr) = a_d p^d + a_{d-1}
p^{d-1} q + \dots + a_0 q^d$ 为整数，且非零
因为 $P$ 无有理根：$\bigl|P\bigl(\frac
pq\bigr)\bigr| \geq q^{-d}$。

\textbf{7.} 注意 $M > 0$：$P'$ 为非零\omterm{def:b1:poly:def}{多项式}（$d \geq
2$），故它不可能在
$\intcc{x-1}{x+1}$ 上恒为零。若 $\bigl|x - \frac pq\bigr| > 1$，则它
显然超过 $\frac{C}{q^d}$。否则 $\frac pq \in
\intcc{x-1}{x+1}$，由中值定理
（\cref{thm:b1:derivative:mvt}）可得 $c$ 介于 $x$ 与
$\frac pq$ 之间，使得
\[
\Bigl|P\Bigl(\frac pq\Bigr)\Bigr|
= \Bigl|P\Bigl(\frac pq\Bigr) - P(x)\Bigr|
= \abs{P'(c)}\,\Bigl|x - \frac pq\Bigr|
\leq M\,\Bigl|x - \frac pq\Bigr| ,
\]
结合第 6 题，$\bigl|x - \frac pq\bigr| \geq
\frac{1}{Mq^d} \geq \frac{C}{q^d}$。

\textbf{8.} 设 $x = \frac ab$ 为刘维尔数。取 $n$ 满足
$2^{n-1} > b$，以及相应的 $\frac pq$，其中 $q \geq 2$：
\[
0 < \Bigl|x - \frac pq\Bigr| < \frac{1}{q^n}
= \frac{1}{q^{n-1}\,q} \leq \frac{1}{2^{n-1} q} <
\frac{1}{bq} ,
\]
与第 1 题矛盾。因此刘维尔数都是无理数。

\textbf{9.} 若刘维尔数 $x$ 是\omterm{pb:b1:logic:1}{代数数}，则它是无理数
（由第 8 题），第 5--7 题给出 $d \geq 2$ 及 $C > 0$，
$\bigl|x - \frac pq\bigr| \geq \frac{C}{q^d}$ 总成立。对
每个 $n$，相应的逼近分数给出 $\frac{C}{q^d} <
q^{-n}$，即 $C < q^{d-n} \leq 2^{d-n}$（因 $q \geq 2$）。
当 $n$ 充分大时，$2^{d-n} < C$，矛盾。刘维尔数
都是\omterm{pb:b1:logic:1}{超越数}。

\textbf{10.} 数字 1 位于第 $1, 2, 6, 24$ 位；前 $25$ 位中的其余
数字皆为 0：
\[
L = 0.1100010000\,0000000000\,00010\dots
\]

\textbf{11.} 当 $m > k$ 时，位置 $n!$（$n > k$）
为两两不同的整数且 $\geq (k+1)!$，故有限等比和
给出
\[
t_m - t_k = \sum_{n=k+1}^{m} 10^{-n!}
\leq \sum_{j = (k+1)!}^{m!} 10^{-j}
< 10^{-(k+1)!}\,\frac{1}{1 - \frac1{10}}
= \frac{10}{9}\,10^{-(k+1)!} ;
\]
对 $m$ 取\omterm{def:b1:reals:bounds}{上确界}，得 $L - t_k \leq
\frac{10}{9}10^{-(k+1)!} < 2\cdot10^{-(k+1)!}$。下界：$L
\geq t_{k+1} = t_k + 10^{-(k+1)!}$。

\textbf{12.} 令 $p_k = 10^{k!}\,t_k \in \N$，$q_k = 10^{k!}$，且
$(k+1)! = (k+1)\,k!$ 可得 $10^{-(k+1)!} = q_k^{-(k+1)}$，
由第 11 题可写为
\[
0 < L - \frac{p_k}{q_k} < \frac{2}{q_k^{\,k+1}} .
\]
给定 $n$，对 $k \geq n$，有 $2\,q_k^{-(k+1)} \leq q_k^{-n}$
（事实上 $q_k^{\,k+1-n} \geq q_k \geq 10 > 2$），且 $q_k \geq 2$，
则满足第 8 题的定义。$L$ 为刘维尔数。

\textbf{13.} 由第 9 题，$L$ 为\omterm{pb:b1:logic:1}{超越数}——这是历史上第一个
被证明为超越的数（刘维尔，1844 年）。
数字交叉验证：该数字串中含有无穷多个 1，
相邻间隔 $(k+1)! - k! = k\cdot k! \to \infty$，因此该数字串
不是最终周期的，又 $L \notin \Q$ 可由
\cref{pb:b1:reals:1}（第 18 题）中的周期性判别法得到，两者一致。

\textbf{14.} 在阶乘位置上取数字 $d_n \in \intint{1}{9}$，
第 11 题的尾项界至多放大
$9$ 倍：$0 < L' - t'_k \leq 9\cdot\frac{10}{9}\,10^{-(k+1)!} =
10\,q_k^{-(k+1)}$（正性源于位置
$(k+1)!$ 上的数字非零）。对 $k \geq n$，有 $10\,q_k^{-(k+1)} \leq
q_k^{-n}$，因 $q_k^{\,k+1-n} \geq 10$，故此数又是刘维尔数，
从而为\omterm{pb:b1:logic:1}{超越数}。这些值两两
互异（对不同的数字选取）；数字串是规范的——
零很多——且规范的数字串决定其值，
见 \cref{pb:b1:reals:1}（第 10 题）。设有任一列表 $k \mapsto
x_k$ 的各数，取第 $k$ 个阶乘位上的数字，
取 $\intint{1}{9}$ 中与 $x_k$ 的该位不同的数字，
所得数即为列表中缺少的同型数。由此可得不可数多个
显式\omterm{pb:b1:logic:1}{超越数}。

\textbf{15.} 先证一个引理：\emph{若 $\bigl|x - \frac pq\bigr|
\geq \frac{C}{q^s}$ 对一切 $\frac pq \neq x$ 成立，则 $x$ 不能
以任何 $\mu > s$ 阶逼近。}事实上，若存在无穷多个
$\frac pq \neq x$ 满足 $\bigl|x - \frac pq\bigr| <
\frac{c}{q^\mu}$，则必迫使 $\frac{C}{q^s} < \frac{c}{q^\mu}$，
即 $q^{\mu - s} < \frac cC$：$q$ 有界，且
距离 $1$ 之内（即与 $x$ 相距）的分数也只有有限多个——
候选者只有有限个，而非无穷多个。现看这一层次：
层次：有理数可以 $1$ 阶逼近（取 $\frac
pq$，$p = \lfloor qx\rfloor + 1$，误差 $\leq \frac1q <
\frac2q$）且不能以任何 $\mu > 1$ 阶逼近（第 1 题给出
引理的前提，取 $s = 1$，$C = \frac1b$）；$\sqrt2$：
可以 $2$ 阶逼近（第 3 题），不能更高（第 2 题及引理）；
每个无理数：至少可以 $2$ 阶逼近（第 4 题）；次数为
$d$ 的\omterm{pb:b1:logic:1}{代数数}：至多可以 $d$ 阶逼近（第 7 题及引理）；刘维尔数：
任意阶（第 12 题中的式子，取 $c = 2$）。

\textbf{16.} 取 $r = \frac ab$，$\frac{p_k}{q_k} + \frac ab =
\frac{b p_k + a q_k}{b q_k} =: \frac{P_k}{Q_k}$, $Q_k = b q_k
\geq 2$，且
\[
\Bigl|(L + r) - \frac{P_k}{Q_k}\Bigr| = L - \frac{p_k}{q_k}
< 2\,q_k^{-(k+1)} = 2\,b^{\,k+1} Q_k^{-(k+1)} .
\]
给定 $n$：对充分大的 $k$，$Q_k^{\,k+1-n} \geq Q_k = b\,10^{k!}
\geq 2\,b^{\,k+1}$（阶乘压过幂），故
误差 $< Q_k^{-n}$，即 $L + r$ 为刘维尔数。由于 $\Q$ \omterm{def:b1:topology:dense}{稠密}
且每个 $L + r$ 为\omterm{pb:b1:logic:1}{超越数}，\omterm{pb:b1:logic:1}{超越数}
在 $\R$ 中\omterm{def:b1:topology:dense}{稠密}。

\textbf{17.} 对 $h \geq 1$，只有有限多个 $P \in
\Z[X]$ 且 $\deg P + \sum_i \abs{a_i} \leq h$（次数 $\leq h$，
且各系数属于 $\intint{-h}{h}$：至多
$(2h+1)^{h+1}$ 个）。每个非零整系数\omterm{def:b1:poly:def}{多项式}都有这样的
高度，且至多有 $\deg P$ 个实根：\omterm{pb:b1:logic:1}{代数数}
构成对 $h$ 的可数并（每项为有限集），因此
可排成单个数列。若\omterm{pb:b1:logic:1}{超越数}也能排成数列，
则将两列交错可列出 $\R$，
与 \cref{pb:b1:reals:1}（第 22 题）矛盾。故
\omterm{pb:b1:logic:1}{超越数}构成不可数集。比较：
康托证明\emph{大多数}实数是\omterm{pb:b1:logic:1}{超越数}，却未举出任何一个；
刘维尔举出了一个，并给出有效常数（第五部分）
——以丰富性证明存在，与以构造证明存在之别。

\textbf{18.} 由第 2 题，引理的前提在 $s
= 2$，$C = \frac14$ 时成立。若 $\sqrt2$ 为刘维尔数，则对 $n =
3$：$\frac{1}{4q^2} < q^{-3}$ 迫使 $q < 4$，故 $q \in \{2,
3\}$；满足 $\frac pq$ 的分母属于这些 $q$ 的分数只有有限多个，
它们中落在 $1$ 之内（以 $\sqrt2$ 为中心）者，各与之相距一个正数
$\geq \varepsilon_0$（$\sqrt2$ 为无理数）；取 $n$ 使
$2^{-n} < \varepsilon_0$，便再无容许的 $\frac pq$：
矛盾。同一论证以 $\frac{C}{q^d}$ 代入，同样表明
任何\omterm{pb:b1:logic:1}{代数数}都不是刘维尔数——
这只是第 9 题换了有效外衣的版本。

\textbf{19.} $99^2 - 2\cdot70^2 = 9801 - 9800 = 1$。于是
\[
\sqrt2 - \frac{99}{70} =
\frac{2 - (99/70)^2}{\sqrt2 + 99/70}
= \frac{-1}{4900\,\bigl(\sqrt2 + \tfrac{99}{70}\bigr)} ,
\qquad
\Bigl|\sqrt2 - \frac{99}{70}\Bigr|
= \frac{1}{4900 \times 2.8284\dots} \approx 7.2\cdot10^{-5} :
\]
$\frac{99}{70} = 1.414285\dots$，对比 $\sqrt2 =
1.414213\dots$——五位数字正确。

\textbf{20.} 对 $P = X^3 - 2$ 作有理根检验：候选根
$\pm1, \pm2, \pm\frac12$，皆非根。故 $d = 3$，第二部分
适用于 $x = 2^{1/3} = 1.2599\dots$。在 $\intcc{x - 1}{x + 1}
\subseteq \intcc{0.25}{2.26}$ 上：$\abs{P'(t)} = 3t^2 \leq 3\,(1 +
2^{1/3})^2 < 3\times(2.26)^2 = 15.32 < 16$，故 $M < 16$，且 $C
\geq \frac{1}{16}$：
\[
\Bigl|2^{1/3} - \frac pq\Bigr| \geq \frac{1}{16\,q^3}
\qquad\text{对一切有理数.}
\]

\textbf{21.} 若 $\bigl|2^{1/3} - \frac pq\bigr| < 10^{-6}$，则
$\frac{1}{16 q^3} < 10^{-6}$，即 $q^3 > \frac{10^6}{16} =
62\,500$；由于 $39^3 = 59\,319 < 62\,500 \leq 64\,000 = 40^3$，
故 $q \geq 40$。

\textbf{22.} 在底 $2$ 下重做第三部分：$B = \sup_k \sum_{n\leq
k} 2^{-n!}$，$q_k = 2^{k!}$，几何尾项（公比
$\frac12$）给出 $0 < B - \frac{p_k}{q_k} \leq 2\cdot2^{-(k+1)!}
= 2\,q_k^{-(k+1)} \leq q_k^{-n}$，对 $k \geq n$。故 $B$ 为
刘维尔数，从而为\omterm{pb:b1:logic:1}{超越数}：论证中的任何部分
都不依赖十进制。

\textbf{23.} 在位置 $3^k$ 处放数字 1：$q_k = 10^{3^k}$，且
尾项界给出 $0 < x^\dagger - \frac{p_k}{q_k} <
2\cdot10^{-3^{k+1}} = 2\,q_k^{-3}$（因 $3^{k+1} = 3\cdot3^k$），
得到无穷多个阶 $3$ 的逼近。由
第 15 题的引理：阶 $3 > 1$ 排除有理性，阶 $3 >
2$ 排除为二次无理数（其刘维尔不等式
中 $s = d = 2$）。但对次数
$\geq 3$ 的\omterm{pb:b1:logic:1}{代数数}，只在阶 $d \geq 3$ 时被排斥：刘维尔的
方法无法将 $x^\dagger$ 与三次\omterm{pb:b1:logic:1}{代数数}分开。这一缺口
由罗斯定理填补——每个代数无理数的逼近阶
恰为 $2$——一个二十世纪的结果，
远超本册范围；承认它，$x^\dagger$ 亦为
\omterm{pb:b1:logic:1}{超越数}。

\textbf{24.} 至多有 $(2H+1)^{d+1}$ 个元组 $(a_0,
\dots, a_d)$，其分量属于 $\intint{-H}{H}$，其中每个非零
\omterm{def:b1:poly:def}{多项式}至多有 $d$ 个实根，故至多产生
$d\,(2H+1)^{d+1}$ 个\omterm{pb:b1:logic:1}{代数数}——这正是
使第 17 题能够将它们全部列出的有限性。

\textbf{25.} (i) 唯一的分析要素是中值定理，
即第 7 题中的定理，将 $P(x) =
0$ 转化为 利普希茨 排斥 $\abs{P(p/q)} \leq M\abs{x -
p/q}$。(ii) 张力在于：整性把 $\abs{P(p/q)}$ 推高到
$q^{-d}$，光滑性则把它压低到 $M\abs{x - p/q}$——若某
有理数离 $x$ 过近，便会被二者夹碎。
(iii) 刘维尔构造出一个带有效常数的\omterm{pb:b1:logic:1}{超越数}；
康托则表明几乎所有实数都是\omterm{pb:b1:logic:1}{超越数}，
却不指名任何一个：构造与\omterm{def:b1:counting:card}{基数}之别。(iv)
\cref{ch:b1:integration}中的周末问题以同一挤压法证明了
$\pi$ 的无理性：一个本应是正整数的积分，
却被困于 $\intoo{0}{1}$ 之中——
以积分代替微分，充当分析的
那一半。
\end{solution}
